\chapter{Short-Horizon Alpha}\label{ch:hf:short-horizon-alpha}

The bid queue is three times the ask queue, and the future ticked up half a second ago. A market maker that ignores both will sell at the ask to
someone who did not. In this chapter's simulated market, where a fund follows its index future half a second late, five features of the book, the
trades and the future forecast the fund's mid two seconds ahead with an information coefficient of 0.48. That is an upper bound: this market's mid
trails its efficient price for seconds, which real books do less. Withdrawing the quote the forecast says will be picked off raises the two-second
mark-out of the market maker's fills from 0.29 to 0.45 ticks a share. The forecast barely ages in one message; in fifty it loses more than a third of
its power.

\section{What predicts the next second}

\begin{definition}[Short-horizon alpha]\label{def:hf:short-horizon-alpha:alpha}
\emph{Short-horizon alpha}\index{short-horizon alpha} is a forecast of an instrument's price change over the next milliseconds to seconds, built
from the state of its order book, its recent trades and the recent moves of related instruments, and used to decide how to quote and when to trade.
\end{definition}

It differs from the intraday alpha of One Quant Book 8, chapter 14 by its horizon and its use. At minutes, a forecast pays for crossing the spread
and is traded like any other signal. At seconds, a forecast is rarely large enough to pay for the spread; its main use is to decide which of the
market maker's own quotes to leave in the book. The predictors are the order-book and trade-flow features of One Quant Book 7, chapters 8 to 10: queue
imbalance, order-flow imbalance, signed trade flow, the lead of a related instrument, and the instrument's own recent move.

\section{Building the forecast: book, flow, cross-asset}

The chapter's market has two instruments on one efficient price. A leads, the index future; B, the fund, follows half a second later, each with its
own order flow. The market maker quotes B and reads A's book. After every message of B, five features are computed over a one-second window
(\cref{lst:hf:short-horizon-alpha:features}): the queue imbalance $(Q^b-Q^a)/(Q^b+Q^a)$; the order-flow imbalance of Cont, Kukanov and Stoikov,
who found price changes over short intervals linear in it with a slope inversely proportional to depth; the net aggressive volume; A's mid change;
and B's own mid change. A ridge regression on four twenty-minute calibration sessions forecasts B's mid change over the next two seconds.

On two validation sessions (information coefficients, signs as fitted, forecasting two seconds ahead):

\begin{center}
\small
\begin{tabular}{@{}lcccccc@{}}
\toprule
& imbalance & order-flow imb. & trade flow & leader's move & own move & combined\\
\midrule
information coefficient & 0.320 & 0.285 & 0.077 & 0.288 & 0.150 & 0.480\\
per unit (ticks) & 0.30 & 0.0045 & 0.0057 & 0.34 & 0.27 & \\
\bottomrule
\end{tabular}
\end{center}

The second row reads: a book fully imbalanced towards the bid forecasts $+0.30$ ticks; a tick of A's move forecasts $+0.34$ ticks of B's; a lot of
order-flow imbalance, $+0.0045$. The own move enters with a positive sign: in this market the mid trends for seconds after the efficient price jumps
(One Quant Book 7 records a lag-one correlation of 0.44 of two-second mid returns in it), and that trend is why every coefficient here is larger than
in real books. Treat the numbers as an upper bound and the ranking as the lesson: the book and the leader carry the forecast, the trade flow little.

\section{Horizon, decay and evaluation at high frequency}

The information coefficient depends on the horizon (\cref{fig:hf:short-horizon-alpha:ic}): 0.28 at half a second, 0.48 at two, a peak of 0.54 at
five, 0.50 at ten. Imbalance peaks early and fades; order-flow imbalance keeps gaining to ten seconds. Kolm, Turiel and Westray, forecasting 115
Nasdaq stocks from order flow with neural networks, found the effective horizon of such forecasts to be about two average price changes; in this
market two price changes take a few seconds, which is where the curve flattens.

\begin{omfigure}
\begin{tikzpicture}
\begin{axis}[width=11cm,height=5.4cm,xlabel={\small forecast horizon (seconds)},ylabel={\small information coefficient},xmode=log,log basis x=10,
  xmin=0.4,xmax=12,xtick={0.5,1,2,5,10},xticklabels={0.5,1,2,5,10},ymin=0,ymax=0.6,tick label style={font=\footnotesize},grid=major,
  grid style={gray!25},table/col sep=comma,legend style={font=\scriptsize,at={(0.5,-0.3)},anchor=north},legend columns=3]
\addplot[black,very thick,mark=*,mark size=1.3pt] table[x=h,y=combined]{figdata/hft/07-short-horizon-alpha/ic.csv};
\addplot[omDef,thick,mark=square*,mark size=1.1pt] table[x=h,y=imbalance]{figdata/hft/07-short-horizon-alpha/ic.csv};
\addplot[omThm,thick,dashed,mark=triangle*,mark size=1.3pt] table[x=h,y=ofi]{figdata/hft/07-short-horizon-alpha/ic.csv};
\addplot[omProp,thick,mark=diamond*,mark size=1.3pt] table[x=h,y=lead]{figdata/hft/07-short-horizon-alpha/ic.csv};
\addplot[gray,thick,dotted,mark=o,mark size=1.1pt] table[x=h,y=own]{figdata/hft/07-short-horizon-alpha/ic.csv};
\addplot[gray!60,thick,mark=x,mark size=1.3pt] table[x=h,y=flow]{figdata/hft/07-short-horizon-alpha/ic.csv};
\legend{combined,imbalance,order-flow imbalance,leader's move,own move,trade flow}
\end{axis}
\end{tikzpicture}

\omcaption{Out-of-sample information coefficient of each feature and of the ridge combination, forecasting the fund's mid change over each
horizon (log scale), on two validation sessions of twenty minutes; the model for each horizon is fitted on four other sessions. Data:
\texttt{hf\_alpha.ic\_by\_horizon}.}\label{fig:hf:short-horizon-alpha:ic}
\end{omfigure}

A forecast is also a race against its own age. The same two-second forecast built from features that are $k$ messages old keeps an information
coefficient of 0.477 one message late (a median of 33 milliseconds here), 0.405 twenty messages late (0.9 seconds) and 0.303 fifty messages late
(2.3 seconds). At the two-second horizon the decay is slow; what decays fast is the value of acting before the others, which is chapter 9's subject.

\section{Coupling the forecast to quotes: skew}

\begin{definition}[Forecast skew]\label{def:hf:short-horizon-alpha:skew}
A \emph{forecast skew}\index{forecast skew} shifts a market maker's quotes in the direction of its short-horizon forecast: the side that the
forecast says will be picked off (the ask when the price is expected to rise) is moved away from the \omterm{def:hf:fair-value:fair}{fair price} or withdrawn, the other side kept or
improved.
\end{definition}

In a small-tick market the skew is continuous: the forecast drift enters the control problem (chapter 4, \omterm{def:hf:extensions-of-the-inventory-framework:drift}{drift-adjusted quoting}), or the \omterm{def:hf:fair-value:fair}{fair price}
itself moves (chapter 2). In a one-tick market it is a decision about which side to show. Cartea, Donnelly and Jaimungal put a volume-imbalance
signal inside a market maker's control problem, calibrated it on eleven Nasdaq stocks over one half-year and tested it over the next, and found that
the signal improved the strategy out of sample.

\begin{omfigure}
\begin{tikzpicture}[font=\small,>=stealth]
\draw[->] (-6,0) -- (6.3,0);
\node[font=\footnotesize] at (0,-1.15) {forecast $\hat\alpha$ (ticks)};
\foreach \x/\l in {-4.5/{$-0.9$},-1.5/{$-0.3$},0/{0},1.5/{$0.3$},4.5/{$0.9$}} {\draw (\x,0.1) -- (\x,-0.1) node[below=1pt,font=\footnotesize] {\l};}
\fill[omDef!15] (-1.5,0.25) rectangle (1.5,1.3);
\node[font=\footnotesize] at (0,0.77) {both sides};
\fill[omThm!15] (1.5,0.25) rectangle (4.5,1.3);
\node[font=\footnotesize] at (3,0.77) {bid only};
\fill[omThm!30] (4.5,0.25) rectangle (6.1,1.3);
\node[font=\footnotesize,align=center] at (5.3,0.77) {bid,\\ and buy};
\fill[omProp!15] (-4.5,0.25) rectangle (-1.5,1.3);
\node[font=\footnotesize] at (-3,0.77) {ask only};
\fill[omProp!30] (-6,0.25) rectangle (-4.5,1.3);
\node[font=\footnotesize,align=center] at (-5.25,0.77) {ask,\\ and sell};
\node[font=\scriptsize,anchor=north] at (3,-0.55) {skew};
\node[font=\scriptsize,anchor=north] at (5.3,-0.55) {take threshold};
\end{tikzpicture}

\omcaption{The chapter's coupling of a two-second forecast to one-tick quotes: both sides are shown while the forecast is within $\pm0.3$ ticks;
beyond it, the side that would be picked off is withdrawn (skew); beyond $\pm0.9$ ticks, half the spread plus the taker fee plus a margin, the
market maker also crosses the spread for one lot (\omterm{def:hf:short-horizon-alpha:take}{take threshold}).}\label{fig:hf:short-horizon-alpha:coupling}
\end{omfigure}

\section{When to take instead of make}

\begin{definition}[Take threshold]\label{def:hf:short-horizon-alpha:take}
A \emph{take threshold}\index{take threshold} is the size of short-horizon forecast beyond which a market maker crosses the spread instead of
waiting to be filled: at least half the spread plus the taker fee, plus a margin for the forecast's error and for the other traders who see the same
signal.
\end{definition}

A market maker that buys at the ask pays half the spread and the taker fee against the mid; it gains only if the mid rises by more before the forecast
is stale. Here half the spread is half a tick, the taker fee a tenth of a tick, the margin three tenths: 0.9 ticks. Forecasts that large are rare:
on the validation sessions 0.01\% of the forecasts exceed it in absolute value, against 5.8\% for the skew threshold of 0.3.

On six test sessions, a one-lot quoter in B (a taker fee of 0.1 cent a share, no rebate):

\begin{center}
\footnotesize
\begin{tabular}{@{}lrrrrr@{}}
\toprule
coupling & passive shares & passive mark-out & takes & take mark-out & P\&L (sd)\\
& a session & (ticks a share) & a session & (ticks a share) & (\$ a session)\\
\midrule
none & 16\,550 & 0.294 & & & $-33.58$ (107.68)\\
skew & 11\,050 & 0.451 & & & 25.58 (50.82)\\
skew and take & 10\,867 & 0.452 & 4.2 & 0.540 & 46.17 (51.22)\\
the same, one message late & 10\,883 & 0.414 & 4.5 & 0.722 & 46.47 (36.29)\\
\bottomrule
\end{tabular}
\end{center}

The skew lifts the passive mark-out by 0.16 ticks a share and gives up a third of the volume, the same trade-off as chapter 6's fading, from a
forecast of direction instead of a forecast of mark-out. Takes are few, four a session; their mark-out is positive but rests on 25 to 27 trades, and
the late version's higher figure is noise. One message late, the passive mark-out drops by 8\%. The P\&L column is too noisy over six sessions to rank
the couplings, as in the previous chapters.

\section{Strategy files}

\begin{strategyfile}{Book-imbalance quote skew}\label{strat:hf:short-horizon-alpha:imbalance}
\sfield{Who pays you, and why} Uninformed takers who cross into the side of the book that is about to hold; informed takers are avoided rather
than paid.
\sfield{Instruments and venues} Large-tick equities, futures and exchange-traded funds on price-time venues.
\sfield{Signal} Queue imbalance $(Q^b-Q^a)/(Q^b+Q^a)$ at the best prices, with order-flow imbalance over a short window; forecast $\hat\alpha$ of the mid
change over one to a few seconds.
\sfield{Sizing and execution} Show both sides while $|\hat\alpha|$ is small; withdraw the side against the forecast beyond a threshold; hysteresis to
keep queue places (chapter 6).
\sfield{Costs} Queue places lost on each withdrawal; messages; the volume given up.
\sfield{How it dies} Every market maker reads the same imbalance, so the side about to hold is crowded and the queue that remains is the one about
to be eaten; no specific public episode is recorded.
\sfield{Horizon, capacity, infrastructure} Seconds; one lot to a few per level; an order-by-order feed and a fast quote update.
\sfield{Backtest honestly} A queue-position model (One Quant Book 7, chapter 18) or a reactive simulator: shadow orders that never lose their place
flatter the withdrawal.
\sfield{Sources} Cartea, Donnelly and Jaimungal (2018); Kolm, Turiel and Westray (2023); this chapter: passive mark-out 0.451 against 0.294 ticks a
share, a synthetic market (upper bound).
\end{strategyfile}

\begin{strategyfile}{Order-flow-imbalance taking}\label{strat:hf:short-horizon-alpha:ofi}
\sfield{Who pays you, and why} Liquidity providers whose quotes are stale after a burst of order flow.
\sfield{Instruments and venues} Liquid futures and equities with deep books.
\sfield{Signal} Order-flow imbalance over the last second, scaled by depth (Cont, Kukanov and Stoikov's linear relation).
\sfield{Sizing and execution} Cross for one lot when $|\hat\alpha|$ exceeds the \omterm{def:hf:short-horizon-alpha:take}{take threshold}; a cooldown between takes.
\sfield{Costs} Half the spread and the taker fee on every trade, against the forecast's gain.
\sfield{How it dies} The same signal is visible to every fast trader; the first to act takes the opportunity, and the arms race shortens its life
(chapter 9).
\sfield{Horizon, capacity, infrastructure} Sub-second to seconds; very small capacity per opportunity; co-location.
\sfield{Backtest honestly} Execute a message later than the signal and at the price then available; count only takes whose forecast exceeded the
threshold at the delayed time.
\sfield{Sources} Cont, Kukanov and Stoikov (2014); this chapter: 4.2 takes a session, mark-out 0.540 ticks on too few trades to trust.
\end{strategyfile}

\begin{strategyfile}{Trade-sign continuation}\label{strat:hf:short-horizon-alpha:sign}
\sfield{Who pays you, and why} The market maker avoids the side a metaorder is taking; the metaorder's owner pays through the spread it keeps
paying.
\sfield{Instruments and venues} Any instrument with long memory in order signs.
\sfield{Signal} Net aggressive volume over a window; the sign autocorrelation of One Quant Book 7, chapter 9.
\sfield{Sizing and execution} A term in the forecast; alone, the weakest of the five features here (information coefficient 0.077).
\sfield{Costs} As for the skew.
\sfield{How it dies} Execution algorithms that randomise sizes and timing hide the sign; persistence is partly offset by mean-reverting limit orders
(Bouchaud and co-authors).
\sfield{Horizon, capacity, infrastructure} Seconds to minutes.
\sfield{Backtest honestly} Sign trades with the rule the live system will use (One Quant Book 7, chapter 9), not with the simulator's truth.
\sfield{Sources} Bouchaud, Gefen, Potters and Wyart (2004); this chapter's coefficient.
\end{strategyfile}

\begin{strategyfile}{Futures-led quoting}\label{strat:hf:short-horizon-alpha:futures}
\sfield{Who pays you, and why} Slower liquidity providers in the fund or the stocks, whose quotes lag the future.
\sfield{Instruments and venues} An index future and its fund or constituents; a Treasury future and the cash bond; any leader and follower.
\sfield{Signal} The leader's mid change over the last window, translated into the follower's units (chapter 2's \omterm{def:hf:fair-value:cross}{cross-instrument fair price}).
\sfield{Sizing and execution} Withdraw the follower's quote on the side the leader has moved against; take when the gap exceeds the \omterm{def:hf:short-horizon-alpha:take}{take threshold}.
\sfield{Costs} Latency between the two venues; fees on both.
\sfield{How it dies} Speed: the median life of an index future-fund arbitrage fell from 97 milliseconds in 2005 to 7 in 2011 as firms invested in
speed (Budish, Cramton and Shim), while its median profit per opportunity stayed near 0.08 index points.
\sfield{Horizon, capacity, infrastructure} Milliseconds to seconds; a data path from the leader's venue, chapter 8.
\sfield{Backtest honestly} Timestamps of both venues on one clock, with the real propagation delay between them.
\sfield{Sources} Budish, Cramton and Shim (2015); this chapter: the leader's move has an information coefficient of 0.288 at two seconds.
\end{strategyfile}

\begin{strategyfile}{Post-sweep reversion}\label{strat:hf:short-horizon-alpha:sweep}
\sfield{Who pays you, and why} A trader who swept several levels and moved the price past its level after the impact decays.
\sfield{Instruments and venues} Stocks and futures with transient impact.
\sfield{Signal} The instrument's own move over the last window, after a sweep; the propagator of One Quant Book 10, chapter 12 gives the expected
reversion.
\sfield{Sizing and execution} Quote into the reversion: show the side that fades the sweep.
\sfield{Costs} Adverse selection when the sweep was informed.
\sfield{How it dies} When the sweep was information, not liquidity: the price keeps going. In this chapter's market the own move predicts
continuation (coefficient $+0.27$), so there is no reversion to trade.
\sfield{Horizon, capacity, infrastructure} Seconds.
\sfield{Backtest honestly} Separate sweeps by what followed them only out of sample; the planted market here would show the strategy losing.
\sfield{Sources} Bouchaud, Gefen, Potters and Wyart (2004); One Quant Book 10, chapter 12.
\end{strategyfile}

\section{Tutorial: forecasting the fund from the future}\label{tut:hf:short-horizon-alpha:tut}

\begin{tutorial}
\textbf{Goal.} Build a seconds-ahead forecast from book, flow and a leading instrument, evaluate it out of sample, and couple it to a quoter.
\textbf{End state:} the three tables, \cref{fig:hf:short-horizon-alpha:ic,fig:hf:short-horizon-alpha:calib}.
\begin{tutsteps}
\item \textbf{Features.} \texttt{firm.hfalpha.FeatureEngine} updates on every message; the same code runs over stored sessions (\texttt{stream})
and inside the quoter.
\omcode{code/firm/hfalpha/firm_hfalpha.py}{56}{79}{Order-flow imbalance, trade flow, own move and imbalance, updated on every message.}\label{lst:hf:short-horizon-alpha:features}
\item \textbf{Fit and evaluate.} \texttt{hf\_alpha.model(h)} fits a ridge regression on four sessions; \texttt{ics(h)} scores it on two others;
\texttt{ic\_by\_delay()} ages the features.
\item \textbf{Couple.} \texttt{AlphaQuoter} reads the book without its own orders, forecasts, skews and takes.
\omcode{code/firm/hfalpha/firm_hfalpha.py}{145}{166}{The forecast drives which side to show and when to cross.}\label{lst:hf:short-horizon-alpha:quoter}
\item \textbf{Compare} with \texttt{compare()} on six test sessions.
\end{tutsteps}
\textbf{What to change next.} Replace the ridge with gradient boosting (One Quant Book 12, chapter 5) and check that the out-of-sample coefficient
improves; make the leader's lag random from session to session and watch the leader's coefficient.
\end{tutorial}

\begin{omfigure}
\begin{tikzpicture}
\begin{axis}[width=11cm,height=5.2cm,xlabel={\small tenth of forecasts},ylabel={\small ticks over two seconds},xtick={1,2,3,4,5,6,7,8,9,10},
  xmin=0.5,xmax=10.5,ymin=-0.4,ymax=0.4,tick label style={font=\footnotesize},grid=major,grid style={gray!25},table/col sep=comma,
  legend style={font=\scriptsize,at={(0.5,-0.3)},anchor=north},legend columns=2]
\addplot[omDef,thick,mark=*,mark size=1.2pt] table[x=group,y=pred]{figdata/hft/07-short-horizon-alpha/calibration.csv};
\addplot[omThm,thick,only marks,mark=square*,mark size=1.5pt] table[x=group,y=real]{figdata/hft/07-short-horizon-alpha/calibration.csv};
\legend{mean forecast,mean realised change}
\end{axis}
\end{tikzpicture}

\omcaption{The two-second forecast on the validation sessions, sorted into tenths: mean forecast and mean realised change of the fund's mid, in
ticks. Data: \texttt{hf\_alpha.calibration}.}\label{fig:hf:short-horizon-alpha:calib}
\end{omfigure}

\section{Build: the short-horizon alpha engine}\label{bld:hf:short-horizon-alpha:hfalpha}

\begin{build}
\bfield{Purpose} One feature engine for research and trading, a forecast fitted and scored out of sample, and the couplings a market maker uses.
\bfield{Interface} \texttt{FeatureEngine(window, lot, leader).update(t, kind, agg, qty, bid, bid\_qty, ask, ask\_qty)}; \texttt{stream(tape, window,
leader)}; \texttt{target(times, mid, horizon)}; \texttt{Ridge.fit(X, y, lam)}, \texttt{.predict}; \texttt{ic}; \texttt{AlphaQuoter(model, window,
leader, skew, take, lag, size, limit)}.
\bfield{Rules} The engine sees only what a feed has shown; the leader is read as of the same time. Features in the quoter use the book without its
own orders, as the model was fitted. Every score is out of sample.
\bfield{Acceptance tests} \texttt{code/firm/hfalpha/tests/}: order-flow imbalance, trade flow, the leader's and own move by hand on a five-message
feed, including the window emptying; the target by hand; the ridge recovers a planted relation.
\bfield{Stretch} Features on several book levels; a forecast per horizon, blended by the quoter's holding time.
\end{build}

\begin{omsources}
\item R.~Cont, A.~Kukanov and S.~Stoikov, ``The price impact of order book events'', \emph{Journal of Financial Econometrics} 12(1), 2014.
\item P.~N.~Kolm, J.~Turiel and N.~Westray, ``Deep order flow imbalance: extracting alpha at multiple horizons from the limit order book'',
\emph{Mathematical Finance} 33(4), 2023.
\item \'A.~Cartea, R.~Donnelly and S.~Jaimungal, ``Enhancing trading strategies with order book signals'', \emph{Applied Mathematical Finance} 25(1),
2018.
\item E.~Budish, P.~Cramton and J.~Shim, ``The high-frequency trading arms race: frequent batch auctions as a market design response'',
\emph{Quarterly Journal of Economics} 130(4), 2015.
\item J.-P.~Bouchaud, Y.~Gefen, M.~Potters and M.~Wyart, ``Fluctuations and response in financial markets: the subtle nature of random price
changes'', \emph{Quantitative Finance} 4(2), 2004.
\end{omsources}

\section{Exercises}

\begin{exercise}[$\star$]\label{exo:hf:short-horizon-alpha:1}
The bid queue holds 12 lots and the ask 4. With the chapter's coefficient of 0.30 ticks per unit of imbalance, what does the imbalance alone
forecast?
\end{exercise}

\begin{exercise}[$\star$]\label{exo:hf:short-horizon-alpha:2}
With a one-cent tick, a half-spread of half a tick and a taker fee of 0.1 cent a share, what is the smallest forecast at which taking breaks even,
before any margin?
\end{exercise}

\begin{exercise}[$\star$]\label{exo:hf:short-horizon-alpha:3}
The best bid rises from 99 to 100 with 3 lots, while the ask stays at 101 and grows from 5 to 6 lots. What is the order-flow imbalance increment?
\end{exercise}

\begin{exercise}[$\star\star$]\label{exo:hf:short-horizon-alpha:4}
Why does the own-move feature have a positive coefficient in this market, and what would you expect in a real stock?
\end{exercise}

\begin{exercise}[$\star\star$]\label{exo:hf:short-horizon-alpha:5}
From \cref{fig:hf:short-horizon-alpha:ic}, which feature is best at half a second and which at ten seconds? What does that suggest about the
coupling for a quoter that holds positions for five seconds?
\end{exercise}

\begin{exercise}[$\star\star$]\label{exo:hf:short-horizon-alpha:6}
Why can a forecast with an information coefficient of 0.48 rarely pay for crossing the spread?
\end{exercise}

\begin{exercise}[$\star\star\star$]\label{exo:hf:short-horizon-alpha:7}
\emph{Coding.} With \texttt{ic\_by\_delay()}, give the information coefficient of the two-second forecast made ten messages late, and the median time
those messages take.
\end{exercise}

\begin{exercise}[$\star\star\star$]\label{exo:hf:short-horizon-alpha:8}
\emph{Find the flaw.} ``Our backtest takes whenever the forecast exceeds the spread and earns 0.7 ticks a trade; it executes at the mid at the
moment of the forecast.''
\end{exercise}

\section{Problem: Three to One on the Bid}

\begin{problem}[{Weekend problem --- three to one on the bid}]\label{pb:hf:short-horizon-alpha:1}
A fund's market maker must decide, message by message, which of its two quotes to show and when to cross.

\medskip\noindent\textbf{Part I --- The forecast.}
\begin{enumerate}
\item Define \omterm{def:hf:short-horizon-alpha:alpha}{short-horizon alpha} and say how it differs from intraday alpha.
\item List the five features and how each is computed.
\item Give the out-of-sample information coefficients at two seconds, and the combination's.
\item Interpret the per-unit coefficients of imbalance and of the leader's move.
\item Why are these coefficients an upper bound for a real market?
\end{enumerate}

\medskip\noindent\textbf{Part II --- Horizons and age.}
\begin{enumerate}[resume]
\item How does the combination's coefficient change from half a second to ten?
\item What did Kolm, Turiel and Westray find about the effective horizon?
\item Give the coefficient one, twenty and fifty messages late, and the times involved.
\item What does \cref{fig:hf:short-horizon-alpha:calib} show?
\end{enumerate}

\medskip\noindent\textbf{Part III --- Using it.}
\begin{enumerate}[resume]
\item Define a \omterm{def:hf:short-horizon-alpha:skew}{forecast skew} and describe it in a one-tick market.
\item Define the \omterm{def:hf:short-horizon-alpha:take}{take threshold} and derive the chapter's 0.9 ticks.
\item How often does the forecast exceed each threshold?
\item Give the couplings' passive shares and mark-outs.
\end{enumerate}

\medskip\noindent\textbf{Part IV --- The verdict.}
\begin{enumerate}[resume]
\item State the \emph{named result}: the passive mark-out per share of the skewed quoter against the unskewed one, and the fraction of it lost one
message late.
\item Why are the take results not evidence?
\item How does the skew relate to chapter 6's fading?
\item Which strategy file has no edge in this market, and why?
\item How did the arms race change futures-led trading between 2005 and 2011?
\item What would you change in the backtest before believing the take mark-out?
\item In one sentence: what is a seconds-ahead forecast mainly for, for a market maker?
\end{enumerate}
\end{problem}

\section{Interview questions}

\begin{interviewq}[$\star$ \iqroles{trader}]\label{iq:hf:short-horizon-alpha:1}
The future just ticked up and your fund quote has not moved. What do you do with your ask, and why?
\end{interviewq}

\begin{interviewq}[$\star\star$ \iqroles{researcher}]\label{iq:hf:short-horizon-alpha:2}
Define order-flow imbalance at the best prices and explain why it predicts price changes better than signed trade volume.
\end{interviewq}

\begin{interviewq}[$\star\star$ \iqroles{researcher}]\label{iq:hf:short-horizon-alpha:3}
How do you evaluate a forecast whose target windows overlap, message after message?
\end{interviewq}

\begin{interviewq}[$\star\star$ \iqroles{developer}]\label{iq:hf:short-horizon-alpha:4}
Research computed features from stored data; production computes them from the live feed. How do you make sure they are the same numbers?
\end{interviewq}

\begin{interviewq}[$\star\star$ \iqroles{trader, researcher}]\label{iq:hf:short-horizon-alpha:5}
When should a market maker take instead of make? Give the threshold and what goes into its margin.
\end{interviewq}

\begin{interviewq}[$\star\star\star$ \iqroles{researcher}]\label{iq:hf:short-horizon-alpha:6}
A forecast $\hat\alpha$ of a price change $y$ has correlation $\rho$ with it, and $\hat\alpha$ is the conditional mean. If $y$ has standard deviation
$s$, how often, roughly, does $|\hat\alpha|$ exceed a threshold $c$ when $\hat\alpha$ is Gaussian? Apply it to $\rho=0.48$, $s=0.6$ ticks, $c=0.9$.
\end{interviewq}
